%%%%%%%%%%%%%%%%%%%%%%%%%%%%%%%%%%%%%%%%%
% University Assignment Title Page 
% LaTeX Template
% Version 1.0 (27/12/12)
%
% This template has been downloaded from:
% http://www.LaTeXTemplates.com
%
% Original author:
% WikiBooks (http://en.wikibooks.org/wiki/LaTeX/Title_Creation)
%
% License:
% CC BY-NC-SA 3.0 (http://creativecommons.org/licenses/by-nc-sa/3.0/)
% 
% Instructions for using this template:
% This title page is capable of being compiled as is. This is not useful for 
% including it in another document. To do this, you have two options: 
%
% 1) Copy/paste everything between \begin{document} and \end{document} 
% starting at \begin{titlepage} and paste this into another LaTeX file where you 
% want your title page.
% OR
% 2) Remove everything outside the \begin{titlepage} and \end{titlepage} and 
% move this file to the same directory as the LaTeX file you wish to add it to. 
% Then add \input{./title_page_1.tex} to your LaTeX file where you want your
% title page.
%
%%%%%%%%%%%%%%%%%%%%%%%%%%%%%%%%%%%%%%%%%
%\title{Title page with logo}
%----------------------------------------------------------------------------------------
%	PACKAGES AND OTHER DOCUMENT CONFIGURATIONS
%----------------------------------------------------------------------------------------

\documentclass[12pt]{article}
\usepackage[english]{babel}
\usepackage[utf8x]{inputenc}
\usepackage{amsmath}
\usepackage{hyperref}
\usepackage[]{algorithm2e}
\usepackage{graphicx}
\usepackage{enumitem}
\usepackage{amssymb}
\usepackage{amsfonts}
\usepackage{multicol}
\usepackage{titlesec}
\usepackage{float}
\usepackage[colorinlistoftodos]{todonotes}
\usepackage{cleveref}



\textheight=250truemm \textwidth=160truemm 
\hoffset=-10truemm \voffset=-20truemm

\begin{document}

\begin{titlepage}

\newcommand{\HRule}{\rule{\linewidth}{0.5mm}} % Defines a new command for the horizontal lines, change thickness here

\center % Center everything on the page
 
%----------------------------------------------------------------------------------------
%	HEADING SECTIONS
%----------------------------------------------------------------------------------------

\textsc{\LARGE Ukrainian Catholic University}\\[1cm] % Name of your university/college
\textsc{\Large  Faculty of Applied Sciences}\\[0.5cm] % Major heading such as course name

%----------------------------------------------------------------------------------------
%	TITLE SECTION
%----------------------------------------------------------------------------------------
\vspace*{1cm}

\HRule \\[0.4cm]
{ \huge \bfseries SLAM Dataset}\\[10pt]
\HRule \\[1cm]
 
%----------------------------------------------------------------------------------------
%	AUTHOR SECTION
%----------------------------------------------------------------------------------------
\vspace*{1cm}

\Large \emph{Authors:}\\
Andrii \textsc{Kryvyi}\\Hordii \textsc{Yelisieiev}\\Oleksandr \textsc{Kosovan}\\Yaroslav \textsc{Prytula}\\[1cm] % Your name

%----------------------------------------------------------------------------------------
%	DATE SECTION
%----------------------------------------------------------------------------------------
\vspace*{1cm}
{\large 2026}\\[2cm] % Date, change the \today to a set date if you want to be precise

%----------------------------------------------------------------------------------------
%	LOGO SECTION
%----------------------------------------------------------------------------------------

\includegraphics[height=5cm]{UCU-Apps.png}\\[1cm] % Include a department/university logo - this will require the graphicx package
 
%----------------------------------------------------------------------------------------

\vfill % Fill the rest of the page with whitespace

\end{titlepage}




\section{Sensor setup}

\subsection{Platform}

For data collection, Clearpath Robotics Husky A200 was used as a base platform. The platform is shown on \cref{fig:husky}.

\begin{figure}
    \centering
    \includegraphics[width=0.5\linewidth]{a200.png}
    \caption{Husky A200 illustration from Husky A200 User Manual}
    \label{fig:husky}
\end{figure}

\subsection{Sensors}

\begin{itemize}
    \item \textbf{\textit{Intel RealSense D435}} --- RGBD Camera. Depth: FOV $\approx 85.2^\circ\times 58^\circ$, $640\times 480$ @ $20$ fps; operating range $\approx 0.3–3$ m. RGB: FOV $\approx 69^\circ\times 42^\circ$, $640\times 480$ @ $20$ fps;

    \item \textbf{\textit{Phidgets Spatial 1042}} --- IMU Sensor. Accelerometer: $\pm 8\text{g}$ @ $200 \text{Hz}$; Gyroscope: $\pm 2000^\circ/\text{s}$ @ $200 \text{Hz}$; Magnetometer: $5.5\text{G}$ @ $125 \text{Hz}$;

    \item \textbf{\textit{Husky Wheel Encoders}} --- x4 Wheel Encoders. Resolution: $78,000\ \text{ticks}/\text{m}$ ($\approx 0.0128\ \text{mm}/\text{tick}$). With $330\ \text{mm}$ wheels, or $\approx 80,865\ \text{ticks}/\text{revolution}$;

    \item  \textbf{\textit{}} --- GNSS Receiver.
\end{itemize}

\section{Calibration} \label{calibration}

\section{Ground truth} \label{groundtruth}

For now, ground truth is estimated solely from GNSS data. This gives only meter-precision and no orientation. The pose is estimated by just linearizing GNSS from Geodetic coordinates to ENU coordinates with fixed origin. For details see \cref{linearization}.

\subsection{GNSS linearization} \label{linearization}

\noindent Since our GNSS sensor provides data as $(P_G=(P_\varphi, P_\lambda, P_h), \Sigma_P)$, where $P_G$ is point in Geodetic coordinates and $\Sigma_P$ is covariance matrix in ENU coordinates with origin $P$. However, we want to have 3d pose with metric axis. We have chosen ENU coordinates for this task, becuase it provides both metric and easy-to-debug system.

\vspace{0.2cm}
\noindent ENU coordinates are defined by constructing a tangent plane to the WGS84 ellipsoid at the given point, which is called an origin $O_G = (O_\varphi, O_\lambda, O_h)$. At the point $O$ we construct three vectors, one pointing to positive latitude (i. e. north-pointing), one pointing to positive longitude (i.e. east-pointing) and one normal to the tangent plane outwards from the center of ellipsoid (i.e. up-pointing). These vectors are normalized to have length 1 meter and are chosen to be basis vectors. To understand better - see the \cref{fig:enu}.

\begin{figure}
    \centering
    \includegraphics[width=0.4\linewidth]{ENU.pdf}
    \caption{ENU(green), ECEF(blue) and Geodetic(yellow) coordinate systems}
    \label{fig:enu}
\end{figure} %% WIKIPEDIA IMAGE

\vspace{0.2cm}
\noindent But before ENU we need to convert Geodetic coordinates $P_G=(P_\varphi, P_\lambda, P_h)$ to ECEF coordinates $P_E = (P_x, P_y, P_z)$ as shown on \cref{fig:enu}. This is done using formulas defining WGS84 surface:

\begin{align*}
    N &= \frac{A}{\sqrt{1 - E^2\sin^2\varphi}}\\
    P_x &= (N + h)\cos \varphi \cos \lambda \\
    P_y &= (N + h)\cos \varphi \sin \lambda \\
    P_z &= (N \cdot (1 - E^2) + h)\sin \varphi \\
\end{align*}

\vspace{0.2cm}
\noindent Where constants are defined as follows:

\begin{align*}
    A &:= 6378137.0 &\text{- }&\text{semi-major axis in meters}\\
    f &:= \frac{1}{298.257223563} &\text{- }&\text{flattening factor} \\
    E^2 &:= f \cdot (2 - f)
\end{align*}

\vspace{0.2cm}
\noindent ECEF is already a metric system, where $P_x$, $P_y$ and $P_z$ are measured in meters. This makes ECEF-to-ENU conversion a simple affine transformation.

\vspace{0.2cm}
\noindent Using multivariate calculus one can derive normalized east, north and up pointing vectors at the point $O$ and construct rotation matrix:

\begin{gather*}
    R(O) = \begin{bmatrix}
        -\sin O_\lambda & \cos O_\lambda & 0 \\
        -\sin O_\varphi \cos O_\lambda & -\sin O_\varphi \sin O_\lambda & \cos O_\varphi \\
        \cos O_\varphi \cos O_\lambda & \cos O_\varphi \sin O_\lambda & \sin O_\varphi
    \end{bmatrix}
\end{gather*}

\begin{figure}
    \centering
    \includegraphics[width=0.8\linewidth]{origin.pdf}
    \caption{Origin location}
    \label{fig:origin}
\end{figure}

\vspace{0.2cm}
\noindent The origin point was chosen arbitrarily at the center of the Campus, and altitude was chosen to match the WGS84 altitude of the ground. See \cref{fig:origin}:

\begin{gather*}
    O_\varphi=49.81754 \qquad O_\lambda=24.02295 \qquad O_h=378.3
\end{gather*}

\vspace{0.2cm}
\noindent Now we can convert our ECEF coordinates to ENU using origin in ECEF and rotation matrix defined before. Also we convert covariance matrix $\Sigma_P$, because it has it's values in ENU coordinates with origin $P$, but we need ENU coordinates with origin $O$. We do it by converting back to ECEF using $R(P)^T$ and then converting to ENU using $R(O)$ via similarity transformations:

\begin{gather*}
    P_T = R(O)\cdot (P_E - O_E) \\
    \Sigma_O =R(O) \cdot R(P)^T \cdot \Sigma_P \cdot R(P) \cdot R(O)^T
\end{gather*}

\vspace{0.2cm}
\noindent In that way we get $P_T = (P_e, P_n, P_u)$, east, north and up components respectively. These are used as $x, y, z$ coordinates for pose. Similarly covariance matrix $\Sigma_O$ measures noise covariance between east, nort and up components respectively.

\subsection{Future improvements}

We plan in future to improve pose by fusing Visual-Inertial odometry with Wheel and GNSS odometry. This might allow us to not only improve pose estimates, but also to estimate orientation at the same time, thus giving us full ground truth information required to evaluate SLAM approaches.

\section{Delivery} \label{delivery}

The dataset is delivered in two formats.

\subsection{Rosbag}

This format uses native ROS2 rosbag via binary messages stored in \texttt{mcap} file, along with metadata. The following topics are present:

\begin{itemize}
    \item \path{/cam_depth/image_raw} --- Intel RealSense D435 raw depth recording (\path{sensor_msgs/Image});
    \item \path{/cam_depth/image_aligned_raw} --- Intel RealSense D435 depth recording aligned to color recording by camera driver (\path{sensor_msgs/Image});
    \item \path{/cam_color/image_raw} --- Intel RealSense D435 raw color recording (\path{sensor_msgs/Image});
    \item \path{/imu} --- Phidgets Spatial 1042 measurements (\path{sensor_msgs/Imu});
    \item \path{/gps} --- SENSOR measurements (\path{sensor_msgs/NavSatFix});
    \item \path{/odom} --- odometry based on Husky Wheel Encoders (\path{nav_msgs/Odometry});
    \item \path{/pose} --- ground truth pose, estimated as described in \cref{groundtruth} (\path{geometry_msgs/PoseWithCovarianceStamped});
    \item \path{/tf} --- ROS TF2 transformations between frames, calibrated as described in \cref{calibration} (\path{tf/tfMessage}).
\end{itemize}

\subsection{TUM-compatible format}

\section{Sequences}

\section{SLAM evaluation}

\begin{thebibliography}{9}


\end{thebibliography}

\end{document}